\begin{table}[htbp]
\centering
\small
\begin{tabularx}{\textwidth}{YL{3.0cm}L{3.0cm}}
\toprule
\textbf{Measure} & \textbf{Training} & \textbf{Test} \\
\midrule
Period & 2 Mar -- 11 Apr 2026 & 13 -- 30 Apr 2026 \\
Visits & 59,139 & 32,014 \\
Visits with 1 / 2 / $\geq$3 stops (\%) & 91.3 / 8.6 / 0.2 & 89.2 / 10.7 / 0.1 \\
Delivery / receive / mixed visits (\%) & 38.3 / 48.5 / 13.2 & 44.2 / 40.8 / 14.9 \\
No RMG / one block / several blocks (\%) & 22.5 / 75.8 / 1.7 & 27.8 / 70.1 / 2.1 \\
\midrule
TTT mean / median / P90 (min) & 36.6 / 29 / 64 & 41.2 / 33 / 75 \\
Median TTT delivery / receive / mixed & 23 / 30 / 41 & 27 / 35 / 48 \\
Median TTT one block / several blocks & 29 / 47 & 33 / 50 \\
Mean approach / stops / transfer / exit & 11.0 / 12.1 / 1.1 / 12.4 & 12.7 / 14.4 / 1.7 / 12.4 \\
Mean trucks in the terminal at Gate In & 77 & 95 \\
Delay-peak hours (\%) & 10.4 & 25.7 \\
\bottomrule
\end{tabularx}
\par\smallskip{\footnotesize\raggedright The approach runs from Gate In to the start of the first stop, the exit from the end of the last stop to Gate Out, and the transfer is the time between stops. Gate events are instantaneous, so the four parts add up to the TTT. Delay-peak hours are Gate-In hours with at least 20 visits and a median TTT of at least 37 minutes, the 90th percentile of the training hours.\par}
\caption{Visit structure and turnaround time (TTT) of the event log.}
\label{tab:obs_descriptives}
\end{table}
